\documentclass[12pt, a4paper]{article}
\usepackage{amsmath, amssymb, amsfonts}
\usepackage{graphicx}
\usepackage{hyperref}
\usepackage[left=2cm, right=2cm, top=2.5cm, bottom=2.5cm]{geometry}
\usepackage{tcolorbox}

\title{Comprehensive Lecture Summaries\\ \Large Introduction to Computational Physics (77315)}
\author{}
\date{Spring Semester, 2026}

\begin{document}
\maketitle
\tableofcontents
\newpage

\section{Lecture 1: 26/03/2026 - Introduction, Roundoff Errors, Numerical Differentiation, and Interpolation}
\label{sec:lec1}

\subsection{Introduction to Computational Physics}
\label{subsec:intro}
Computational physics combines numerical analysis methods to solve physical problems. Historically, physics progressed through a binary of theoretical physics (creating mathematical models) and experimental physics (measuring reality). Today, the computational approach serves as a formidable third branch. By leveraging the immense processing power of computers, we can solve complex models (such as nonlinear differential equations with no analytical solutions) and perform macroscopic simulations of many-body systems that would be otherwise impossible to track.

When translating pure mathematics into computer code, we must abandon the notion of "infinite precision." Three main factors must always be considered:
\begin{enumerate}
    \item \textbf{Roundoff Errors}: Inaccuracies due to the finite memory computers use to store numbers.
    \item \textbf{Truncation Errors (Accuracy)}: Inaccuracies arising from using a finite mathematical approximation (like stopping a Taylor series early).
    \item \textbf{Stability}: The tendency of an algorithm to either dampen or exponentially amplify these microscopic errors over time.
\end{enumerate}

\subsection{Errors, Accuracy, and Stability}
\label{subsec:errors}

\subsubsection{Catastrophic Cancellation and Roundoff Errors}
Computers represent real numbers using a limited number of bits. A common representation is 4 bytes (Float32, yielding $\sim 7$ decimal digits of precision) or 8 bytes (Float64, yielding $\sim 16$ decimal digits). 

Because of this finite representation, every basic arithmetic operation introduces a tiny relative error. While this is usually negligible, a severe problem known as \textbf{catastrophic cancellation} occurs when subtracting two very close numbers. 

\begin{tcolorbox}[colback=blue!5!white,colframe=blue!75!black,title=Pedagogical Example: The Quadratic Formula]
Consider the quadratic equation $x^2 - 1000x + 0.001 = 0$. Here, $b = 500$ and $c = 0.001$.
The standard quadratic formula gives the roots as $x = b \pm \sqrt{b^2 - c}$. 

If we calculate the smaller root using subtraction:
$x_{small} = 500 - \sqrt{250000 - 0.001} \approx 500 - 499.999999 = 0.000001$.

If our computer only has 7 digits of precision (Float32), $\sqrt{250000 - 0.001}$ will be rounded to exactly $500.0000$. The subtraction $500.0000 - 500.0000$ yields $0$. We have completely lost the true answer due to catastrophic cancellation!

\textbf{The elegant solution}: Multiply the numerator and denominator by the conjugate $b + \sqrt{b^2-c}$. This transforms the formula to $x_{small} = \frac{c}{b + \sqrt{b^2 - c}}$.
Now: $x_{small} = \frac{0.001}{500 + 500} = \frac{0.001}{1000} = 10^{-6}$.
By avoiding the subtraction of close numbers, we maintain perfect precision.
\end{tcolorbox}

\subsubsection{Truncation Error and Stability}
Truncation error is the difference between the exact mathematical formula and the approximated algorithm. For example, a Taylor series is infinite, but a computer can only sum $N$ terms. The discarded terms constitute the truncation error.

\textbf{Stability} dictates whether an algorithm is usable. An algorithm is unstable if tiny roundoff or truncation errors grow exponentially as the computation proceeds.

\begin{tcolorbox}[colback=red!5!white,colframe=red!75!black,title=Intuitive Analogy: The Golden Mean Recursion]
Consider calculating the powers of the "Golden Mean" $\phi = \frac{\sqrt{5}-1}{2} \approx 0.618$.
Mathematically, $\phi$ satisfies the recurrence relation: $\phi^{n+1} = \phi^{n-1} - \phi^n$.
This seems like a brilliantly fast way to compute $\phi^{100}$ without multiplication!

However, this linear recurrence relation has two mathematical roots: $\phi_1 \approx 0.618$ and $\phi_2 = \frac{-\sqrt{5}-1}{2} \approx -1.618$.
Even if we start with the exact values for $n=1, 2$, the tiny roundoff error introduced by the computer acts as a "seed" for the second root. Because $|\phi_2| > 1$, its contribution grows exponentially. By $n=50$, the parasitic $\phi_2$ term completely dominates, and the algorithm blows up to infinity. This is the hallmark of numerical instability.
\end{tcolorbox}

\subsection{Numerical Differentiation}
\label{subsec:diff}
How do we teach a computer to calculate a derivative, $\lim_{h \to 0} \frac{f(x+h)-f(x)}{h}$? We cannot set $h=0$, nor can we make $h$ infinitely small due to roundoff errors. 

If we have a discrete grid of points $x_n = x_0 + nh$, we can approximate derivatives using Taylor expansions:
\begin{align*}
f(x_0 + h) &= f_0 + h f' + \frac{h^2}{2} f'' + \frac{h^3}{6} f''' + \mathcal{O}(h^4) \\
f(x_0 - h) &= f_0 - h f' + \frac{h^2}{2} f'' - \frac{h^3}{6} f''' + \mathcal{O}(h^4)
\end{align*}
Subtracting the second equation from the first gracefully cancels out the $f''$ term, yielding the \textbf{Central Difference Formula}:
\begin{equation}
f' = \frac{f_{+1} - f_{-1}}{2h} + \mathcal{O}(h^2)
\end{equation}
Geometrically, instead of drawing a secant line from $x$ to $x+h$ (which is biased), we draw a secant line from $x-h$ to $x+h$, which perfectly balances the slope around $x$.

\textbf{The Step-Size Dilemma ($h$)}: If we make $h$ large, the Taylor approximation fails (huge truncation error $e_t \sim h^2 f'''$). If we make $h$ microscopic, catastrophic cancellation destroys the numerator $f_{+1} - f_{-1}$ (huge roundoff error $e_r \sim \frac{\epsilon_f f}{h}$). 
To find the optimal $h$, we minimize the total error $E(h) = \frac{\epsilon_f f}{h} + h^2 f'''$. Taking the derivative with respect to $h$ and equating to zero yields $h_{opt} \sim \sqrt[3]{\frac{3 \epsilon_f f}{f'''}}$.

\subsection{Interpolation}
\label{subsec:interp}
Interpolation is the art of reading "between the lines" of a discrete dataset. 
\begin{itemize}
    \item \textbf{Lagrange Interpolation}: Fits a single global polynomial of degree $N-1$ through exactly $N$ points. While elegant, high-degree polynomials suffer from \textbf{Runge's Phenomenon}: they oscillate violently at the edges of the interval. 
    \item \textbf{Cubic Splines}: To avoid global oscillations, we use piecewise polynomials. Imagine taking a flexible piece of wood (a ``spline'' used by historical draftsmen) and pinning it at the data points. The wood naturally bends to minimize potential energy, creating a smooth curve. Mathematically, a Cubic Spline enforces that the function value, first derivative (slope), and second derivative (curvature) are all continuous across the pins. This local continuity naturally results in a highly stable, tri-diagonal system of equations that is trivial to solve.
\end{itemize}

\begin{tcolorbox}[colback=green!5!white,colframe=green!50!black,title=The Tridiagonal System for Cubic Splines]
Let $h_i = x_{i+1} - x_i$ be the spacing and $m_i = S''(x_i)$ be the unknown second derivative (``moment'') at each knot. Enforcing continuity of $S$, $S'$, and $S''$ across each interior knot yields the tridiagonal recurrence:
\begin{equation}
\label{eq:spline_tridiag}
h_i\, m_{i-1} + 2(h_i + h_{i+1})\, m_i + h_{i+1}\, m_{i+1} = 6\left(\frac{y_{i+1}-y_i}{h_{i+1}} - \frac{y_i - y_{i-1}}{h_i}\right)
\end{equation}
With natural boundary conditions $m_0 = m_N = 0$, this gives an $(N-1)\times(N-1)$ tridiagonal linear system. By the band-matrix optimisation discussed in Section~\ref{subsec:bandmatrices}, this system is solved in $\mathcal{O}(N)$ time — far cheaper than the $\mathcal{O}(N^3)$ of full Gaussian elimination.
\end{tcolorbox}

\vspace{0.5cm}
\noindent \textbf{Further Reading for Lecture \ref{sec:lec1}:}
\begin{itemize}
    \item \textit{Numerical Recipes}, W.H. Press et al. - Chapter 1 (Preliminaries on Errors), Chapter 5.3 (Polynomial Interpolation), Chapter 5.7 (Numerical Derivatives).
    \item \textit{Computational Physics}, Mark Newman - Chapter 4 (Accuracy and Speed), Chapter 5 (Integrals and Derivatives).
\end{itemize}

\newpage
\section{Lecture 2: 12/04/2026 - Numerical Integration, Root Finding, and Linear Equations}
\label{sec:lec2}

\subsection{Numerical Integration (Quadrature)}
\label{subsec:integration}
Numerical integration calculates the area under a curve $\int_a^b f(x)dx$ by evaluating the function at specific discrete points.

\begin{itemize}
    \item \textbf{Simpson's Rule}: Rather than approximating the area using flat rectangles or slanted trapezoids, Simpson's rule groups points in sets of three and fits a parabola through them. The area under this parabola forms the integral estimate:
    \begin{equation}
    \int_{-h}^{h} f(x)dx = \frac{h}{3} [f(-h) + 4f(0) + f(h)]
    \end{equation}
    It is remarkably accurate, possessing a global error of $\mathcal{O}(h^4)$.
    
    \item \textbf{Gauss-Legendre Integration}: This method requires a profound shift in intuition. In Simpson's rule, our $x$-coordinates are fixed on a rigid, evenly spaced grid. Gauss asked: \textit{What if we leave the $x$-coordinates as free variables?} By optimizing both the weights $w_i$ AND the sampling locations $x_i$, we effectively double our degrees of freedom! 
    The optimal $x_i$ coordinates perfectly match the roots of the Legendre polynomials. Using just $n$ optimally placed points, Gauss-Legendre quadrature can perfectly integrate any polynomial up to degree $2n-1$.
    
    \item \textbf{Recursive (Adaptive) Integration}: Imagine integrating a function that is completely flat for $90\%$ of the domain, but oscillates wildly in the last $10\%$. A fixed grid would waste massive computational resources calculating the flat area. Adaptive integration solves this by evaluating the integral on a coarse interval, splitting the interval in half, and re-evaluating. If the two answers agree within a tolerance $\epsilon$, it stops. If they disagree, it recursively dives deeper \textit{only} in the problematic, highly-varying sub-regions.
\end{itemize}

\subsection{Linear Equations (Gaussian Elimination)}
\label{subsec:lineareq}
Solving a system of $n$ linear equations $Ax = b$ is the backbone of almost all computational physics.
The standard algorithm is \textbf{Gaussian Elimination}, which systematically subtracts rows to eliminate variables, transforming $A$ into an upper-triangular matrix.
This sweeping process takes $\mathcal{O}(n^3)$ operations. Once triangular, the system is trivially solved from the bottom up via \textbf{Back Substitution} in $\mathcal{O}(n^2)$ operations.

\begin{tcolorbox}[colback=green!5!white,colframe=green!50!black,title=The Critical Necessity of Partial Pivoting]
During Gaussian elimination, we divide the entire row by the "pivot" element $A_{ii}$. 
Imagine the pivot element is physically non-zero, but incredibly small, say $10^{-10}$. 
Dividing by $10^{-10}$ is equivalent to multiplying the entire row by $10^{10}$! This massive amplification instantly takes microscopic roundoff errors residing in the other columns and blows them up to macroscopic, catastrophic proportions, destroying the rest of the matrix.

\textbf{Partial Pivoting} rescues the algorithm. Before processing a column, we scan down the rows to find the element with the largest absolute magnitude. We physically swap that row to the top. By ensuring we always divide by the largest possible number, we suppress error amplification and guarantee numerical stability.
\end{tcolorbox}

\vspace{0.5cm}
\noindent \textbf{Further Reading for Lecture \ref{sec:lec2}:}
\begin{itemize}
    \item \textit{Numerical Recipes}, W.H. Press et al. - Chapter 4 (Integration of Functions), Chapter 2.2 (Gaussian Elimination with Backsubstitution).
    \item \textit{Computational Physics}, Mark Newman - Chapter 5 (Integrals and Derivatives), Chapter 6 (Solution of Linear and Nonlinear Equations).
\end{itemize}

\newpage
\section{Lecture 3: 23/04/2026 - Band Matrices, Least Squares, Gram-Schmidt, and Eigenvalues}
\label{sec:lec3}

\subsection{Large and Sparse Matrices (Band Matrices)}
\label{subsec:bandmatrices}
In real physical systems, ``everything does not interact with everything.'' Consider the 1D heat diffusion equation on a wire discretised into 10,000 segments. The temperature at segment $i$ only directly influences its immediate neighbours, $i-1$ and $i+1$.
If we build the interaction matrix $A$, it will be $10000 \times 10000$ (100 million elements), but $99.99\%$ of those elements will be zero! The only non-zero elements lie on the main diagonal and the two adjacent diagonals. This is a \textbf{Band Matrix}. (The cubic spline system in Section~\ref{subsec:interp} is another classic example of a naturally occurring tridiagonal system.)

By explicitly telling our computer to only store and operate on the band of width $n_b$, the Gaussian elimination time plummets from $\mathcal{O}(n^3)$ to $\mathcal{O}(n_b^2 \cdot n)$. Memory consumption shrinks from Gigabytes to mere Megabytes.

\subsection{Least Squares}
\label{subsec:leastsquares}
When we fit a line to noisy experimental data, we have an overdetermined system $Ax = b$ with more equations (data points) than unknowns (line parameters). A perfect solution is impossible. 

\textbf{Geometric Intuition}: The columns of $A$ span a vector space (the "model space"). The experimental data vector $b$ lives outside this space due to noise. To find the "best" parameters $x$, we must drop a perpendicular shadow (an orthogonal projection) from $b$ down onto the model space.
The error vector $(b - Ax)$ must be orthogonal to every column of $A$. Therefore, multiplying by $A^T$ yields zero:
$A^T (b - Ax) = 0 \implies A^T A x = A^T b$. 
These are the famous \textbf{Normal Equations}. Solving them minimizes the squared error.

\subsection{Gram-Schmidt Orthogonalization}
\label{subsec:gramschmidt}
Explicitly computing the matrix $A^T A$ (as required by the Normal Equations in Section~\ref{subsec:leastsquares}) is numerically dangerous because squaring the condition number of the matrix violently amplifies floating-point inaccuracies.
Instead, we can use the \textbf{Gram-Schmidt process} to orthogonalise the columns of $A$ directly.

\textbf{Geometric Analogy}: Imagine two non-orthogonal vectors, a ladder leaning against a wall, and the floor. We want to construct a pure vertical vector. We shine a light straight down, find the shadow (projection) of the ladder on the floor, and subtract that shadow from the ladder. What remains is a perfectly vertical vector orthogonal to the floor.
Mathematically, for each vector $a_j$, we subtract its projection onto the previously established orthogonal basis vectors $a_1$:
\begin{equation}
a'_j = a_j - \frac{a_1 \cdot a_j}{a_1 \cdot a_1} a_1
\end{equation}
Repeating for all $j = 2, \ldots, m$ yields a set of mutually orthogonal columns.

\textbf{Connection to QR Decomposition}. The Gram-Schmidt process is algebraically equivalent to a QR factorisation of $A$. Define the accumulated transformation matrix $R$ (upper-triangular) such that $A = \hat{A} R$, where $\hat{A}$ has orthonormal columns ($\hat{A}^T \hat{A} = I$). Substituting into the Normal Equations:
\begin{equation}
A^T A x = A^T b \;\implies\; R^T \underbrace{\hat{A}^T \hat{A}}_{=\, I} R\, x = R^T \hat{A}^T b \;\implies\; Rx = \hat{A}^T b
\end{equation}
The right-hand side $\hat{A}^T b$ is a simple matrix-vector product, and $Rx = \hat{A}^T b$ is an upper-triangular system solved by back-substitution in $\mathcal{O}(m^2)$. This entirely bypasses forming $A^T A$.

\begin{tcolorbox}[colback=orange!5!white,colframe=orange!75!black,title=Numerical Stability Warning: Loss of Orthogonality]
Each Gram-Schmidt step subtracts a projection from a vector. If the columns of $A$ are \textit{nearly linearly dependent}, the remainder after subtraction can be many orders of magnitude smaller than the original. Once intermediate values approach machine epsilon ($\varepsilon_{\text{mach}} \approx 10^{-15}$ for Float64), the subtraction is dominated by floating-point noise and \textbf{orthogonality is silently lost}. The resulting vectors are not truly perpendicular, and the solution degrades without warning. In practice this is addressed by ``Modified Gram-Schmidt'' (subtracting projections one at a time rather than all at once) or by Householder reflections.
\end{tcolorbox}

\subsection{Introduction to Eigenvalues --- Jacobi Method}
\label{subsec:eigenintro}
For a real symmetric matrix $A$, we seek the eigenvalues $\lambda_i$ such that $A\varphi_i = \lambda_i \varphi_i$.
If the matrix were purely diagonal, the eigenvalues would just be the numbers sitting on the diagonal. The \textbf{Jacobi Method} strives to achieve this by applying a sequence of plane-rotation transformations $T^T A T$, where each $T$ is a Givens rotation matrix that zeros one off-diagonal pair $(A_{pq}, A_{qp})$.
Intuitively, at each step we take the 2D ``slice'' of the $N$-dimensional matrix containing the largest off-diagonal element and physically rotate our coordinate system in that plane until that element vanishes. By iteratively rotating, we ``sweep'' the numerical weight from the off-diagonal entries onto the main diagonal. Convergence is guaranteed by two key invariants:

\begin{enumerate}
    \item \textbf{Trace Invariance}: $\mathrm{Tr}(T^T A T) = \mathrm{Tr}(A)$. The sum of the diagonal elements — equal to the sum of all eigenvalues — never changes.
    \item \textbf{Frobenius Norm Invariance}: $\sum_{i,j}(T^T A T)_{ij}^2 = \sum_{i,j} A_{ij}^2$. The total squared ``weight'' of the matrix is conserved under any orthonormal transformation.
\end{enumerate}

\begin{tcolorbox}[colback=green!5!white,colframe=green!50!black,title=Why the Algorithm Must Converge]
The Frobenius norm is shared between diagonal and off-diagonal entries. Since the total is fixed, any weight removed from the off-diagonal must go to the diagonal. Each rotation explicitly sets $A_{pq} = A_{qp} = 0$, transferring $2A_{pq}^2$ of squared weight from $S_{\text{off}} = \sum_{i \neq j} A_{ij}^2$ to the diagonal entries $A_{pp}$ and $A_{qq}$.

Crucially, this transfer is \emph{irreversible in net}: although later rotations may reintroduce small values at $(p,q)$, the \emph{sum} $S_{\text{off}}$ is bounded above by:
\begin{equation}
S_{\text{off}}^{(k+1)} \leq S_{\text{off}}^{(k)} \cdot \left(1 - \frac{2}{n(n-1)}\right)
\end{equation}
after each single rotation. Over a full cycle of $n(n-1)/2$ rotations, $S_{\text{off}}$ decreases by at least a factor of $e^{-1}$, so convergence to a diagonal matrix is guaranteed. (The full convergence analysis is given in Section~\ref{subsec:jacobi_convergence}.)
\end{tcolorbox}

\vspace{0.5cm}
\noindent \textbf{Further Reading for Lecture \ref{sec:lec3}:}
\begin{itemize}
    \item \textit{Numerical Recipes}, W.H. Press et al. - Chapter 2.4 (Band Diagonal Systems), Chapter 15.4 (General Linear Least Squares), Chapter 11.0 (Eigensystems Introduction).
    \item \textit{An Introduction to Computational Physics}, Tao Pang - Chapter 4 (Matrix Operations).
\end{itemize}

\newpage
\section{Lecture 4: 26/04/2026 - Jacobi Convergence, Multidimensional Roots, and Minimization}
\label{sec:lec4}

\subsection{Jacobi Method: Convergence Rate and Complexity}
\label{subsec:jacobi_convergence}
Continuing from the introduction in Section~\ref{subsec:eigenintro}, we now quantify convergence speed and describe an efficient implementation.

\textbf{Convergence Rate per Cycle.} Define a \textbf{cycle} as one complete pass through all $\frac{n(n-1)}{2}$ off-diagonal element pairs. Let $S^{(k)} = \sum_{i \neq j} A_{ij}^2$ after $k$ cycles. Applying the per-rotation bound repeatedly over a full cycle gives:
\begin{equation}
S^{(1)} \leq S^{(0)} \cdot \left(1 - \frac{2}{n(n-1)}\right)^{n(n-1)/2} \approx S^{(0)} \cdot e^{-1}
\end{equation}
Each cycle reduces the off-diagonal weight by a factor of roughly $e^{-1} \approx 0.37$. Consequently:
\begin{itemize}
    \item \textbf{7 cycles} reduce $S$ by $e^7 \approx 10^3$ (one part per thousand).
    \item \textbf{14 cycles} reduce $S$ by $e^{14} \approx 10^6$ (one part per million).
\end{itemize}
In practice, 10 cycles is sufficient for nearly all applications.

\textbf{Computational Complexity.} Each rotation modifies 2 rows and 2 columns of $n$ elements, requiring $\sim 12$ multiplications per zero created. A full cycle has $\frac{n(n-1)}{2}$ rotations; for 10 cycles the total operation count is approximately $60\,n^3$, confirming $\mathcal{O}(n^3)$ scaling.

\textbf{Tracking Eigenvectors.} The eigenvalues appear as the diagonal entries after convergence. To also recover eigenvectors, maintain the accumulated rotation matrix $T = T_1 T_2 \cdots T_k$, initialised as $T = I$. After each plane rotation parametrised by $\cos\phi = c$ and $\sin\phi = s$ acting on the $(p,q)$ plane, update two columns of $T$ for all rows $r$:
\begin{align}
T_{r,p} &\leftarrow c\, T_{r,p} - s\, T_{r,q} \\
T_{r,q} &\leftarrow s\, T_{r,p} + c\, T_{r,q}
\end{align}
At convergence the columns of $T$ are the eigenvectors $\varphi_i$.

\begin{tcolorbox}[colback=blue!5!white,colframe=blue!75!black,title=Optimisation: Row-Maximum Vector ($\mathcal{O}(n^4) \to \mathcal{O}(n^3)$)]
\textbf{The naive bottleneck}: scanning the full upper-triangular part to find $\max_{i<j}|A_{ij}|$ costs $\mathcal{O}(n^2)$ per rotation, making the total $\mathcal{O}(n^4)$.

\textbf{The fix}: maintain a vector $\mathtt{j\_max}[i]$ storing the column index of the largest off-diagonal element in row $i$.
\begin{enumerate}
    \item \textbf{Initialise} ($\mathcal{O}(n^2)$, once): scan each row and fill $\mathtt{j\_max}$.
    \item \textbf{Find global max} ($\mathcal{O}(n)$): compare $A_{i,\,\mathtt{j\_max}[i]}$ over all $i$.
    \item \textbf{After each rotation} (modifies rows and columns $p$ and $q$):
    \begin{itemize}
        \item Re-scan rows $p$ and $q$ to update $\mathtt{j\_max}[p]$ and $\mathtt{j\_max}[q]$ ($\mathcal{O}(n)$).
        \item For every other row $i$: if $\mathtt{j\_max}[i] \in \{p,q\}$, the stored maximum is stale --- re-scan row $i$ ($\mathcal{O}(n)$). This happens with probability $\sim 2/n$, giving expected cost $\mathcal{O}(1)$ per row.
        \item Otherwise, compare $A_{i,\,\mathtt{j\_max}[i]}$ against the two updated values $A_{i,p}$ and $A_{i,q}$ to check for a new maximum ($\mathcal{O}(1)$ per row).
    \end{itemize}
\end{enumerate}
Net result: $\mathcal{O}(n)$ expected cost per rotation $\Rightarrow$ total complexity back to $\mathcal{O}(n^3)$.
\end{tcolorbox}

\subsection{Finding the Root of a Function in Multiple Dimensions}
\label{subsec:multiroot}
In 1D, finding a root is straightforward: if $f(a) < 0$ and $f(b) > 0$, the Intermediate Value Theorem guarantees a root in $(a,b)$, and bisection halves the bracket at every step. In $N$ dimensions, solving $F(x) = 0$ is vastly harder.

\textbf{Why Bracketing Fails in $N$ Dimensions.} Consider just $N=2$: we have two scalar functions $F_1(x,y)$ and $F_2(x,y)$, and we seek points where both vanish simultaneously. The zero-set of $F_1$ is a curve in the $(x,y)$ plane; likewise for $F_2$. The roots are the intersections of these two curves. There is no bounding rectangle whose corner values can certify the presence of an intersection inside --- the curves may bend, avoid each other entirely, or cross multiple times with no detectable sign pattern at the corners. In three or more dimensions the situation is geometrically even less tractable. We therefore have no global ``trapping'' mechanism, and we must resort to local iterative methods that depend sensitively on the starting point. Convergence to a particular root is not guaranteed, and we may miss roots entirely.

\textbf{Multidimensional Newton-Raphson.} Geometrically, if we are standing on an $N$-dimensional curved landscape, we compute the slope in every direction to construct a flat tangent hyperplane, then slide to where that plane intersects zero.
The slopes form the \textbf{Jacobian matrix} $J$ with entries $J_{ij} = \frac{\partial F_i}{\partial x_j}$. The Newton step $\delta$ satisfies:
\begin{equation}
J\,\delta = -F(x)
\end{equation}
This is just a square linear system solved by Gaussian elimination (Section~\ref{subsec:lineareq}). The new estimate is $x_{\mathrm{new}} = x + \delta$, and the process repeats until $|F(x)|$ is acceptably small. If the starting point is sufficiently close to the root, Newton-Raphson converges \textit{quadratically} (the number of correct digits doubles each iteration). However, a poor starting point can send the iterates diverging to infinity --- which motivates the line-search strategy described in Section~\ref{subsec:linesearch}.

\subsection{Line Search and Backtracking}
\label{subsec:linesearch}
The danger of Newton-Raphson is that the flat tangent plane is a terrible approximation of the curved landscape if we take a massive step; it can easily shoot us off into infinity.

\begin{tcolorbox}[colback=yellow!5!white,colframe=yellow!60!black,title=The Blind Hiker Analogy]
Imagine a blindfolded hiker trying to find the lowest point of a valley. Newton's method tells the hiker exactly which direction points downhill, and roughly how far to walk.
However, if the valley is narrow, taking a massive step might cause the hiker to walk past the bottom and ascend high up the opposite mountain face!

To prevent this, we construct a scalar objective function $f(x) = \frac{1}{2} F(x) \cdot F(x)$.
The hiker calculates the full Newton step $p$, but doesn't take it blindly. They take a step of length $\lambda p$. They check their altitude $f(x_{new})$. If their altitude is higher than where they started, they realise they overstepped! They \textbf{backtrack}, dynamically reducing $\lambda$ using parabolic curve fitting to guess a shorter, safer step size, ensuring every single step strictly lowers their altitude.
\end{tcolorbox}

\textbf{Descent Guarantee.} Define the merit function $f(x) = \tfrac{1}{2}F(x)\cdot F(x)$ and the Newton direction $p = -J^{-1}F(x)$. Then:
\begin{equation}
\nabla f \cdot p = F^T J \cdot (-J^{-1}F) = -F \cdot F < 0
\end{equation}
so $p$ is always a \textit{descent direction} for $f$. For a sufficiently small $\lambda > 0$, the step $x + \lambda p$ will lower $f$. We accept $\lambda$ as soon as the \textbf{sufficient-decrease (Armijo) condition} is satisfied:
\begin{equation}
f(x + \lambda p) \leq f(x) + \varepsilon\,\nabla f(x) \cdot (\lambda p), \qquad \varepsilon \approx 10^{-4}
\end{equation}
Starting from $\lambda = 1$ (the full Newton step), we reduce $\lambda$ by fitting a polynomial to the 1D function $g(\lambda) = f(x + \lambda p)$ and jumping to its predicted minimum.

\begin{tcolorbox}[colback=blue!5!white,colframe=blue!75!black,title=Backtracking: Parabolic and Cubic Polynomial Fitting]
\textbf{First backtrack --- parabolic fit.} Three conditions are known: $g(0)$, $g'(0) = \nabla f \cdot p < 0$, and $g(1)$. A parabola $g \approx a\lambda^2 + b\lambda + c$ through these points has its minimum at:
\begin{equation}
\lambda^* = -\frac{g'(0)}{2\bigl[g(1) - g(0) - g'(0)\bigr]}
\end{equation}
Clamp $\lambda^*$ to $[0.1,\, 0.5]$ to prevent steps that are either negligibly small or dangerously large.

\textbf{Subsequent backtracks --- cubic fit.} Four conditions are now available: $g(0)$, $g'(0)$, $g(\lambda_1)$ (first guess), and $g(\lambda_2)$ (second guess). Write $g(\lambda) = a\lambda^3 + b\lambda^2 + g'(0)\lambda + g(0)$ and solve for $(a,b)$ via:
\begin{equation}
\frac{1}{\lambda_1 - \lambda_2}
\begin{pmatrix} \lambda_1^{-2} & -\lambda_2^{-2} \\ -\lambda_1^{-3} & \lambda_2^{-3} \end{pmatrix}
\begin{pmatrix} g(\lambda_1) - g(0) - g'(0)\lambda_1 \\ g(\lambda_2) - g(0) - g'(0)\lambda_2 \end{pmatrix}
= \begin{pmatrix} a \\ b \end{pmatrix}
\end{equation}
The minimum of the cubic is at $\lambda^* = \bigl(-b + \sqrt{b^2 - 3a\,g'(0)}\bigr)/(3a)$, again clamped to a safe interval.

If no $\lambda$ satisfies the Armijo condition, the starting point is outside the basin of attraction and a different initial $x_0$ must be chosen.
\end{tcolorbox}

\subsection{Finding a Minimum in One Dimension - Golden Section Search}
\label{subsec:1dmin}
To find the minimum of a 1D function without calculating derivatives, we "bracket" the minimum using three points $a < b < c$ shaped like a V. We evaluate a new point $x$ inside the larger interval to slowly shrink the bracket.

Where is the mathematically optimal place to put $x$? We want to guarantee that no matter whether the new bracket forms on the left or the right, the fractional reduction in the bracket's size is identical.
This geometric constraint forces the intervals to obey the \textbf{Golden Ratio}:
\begin{equation}
w = \frac{3 - \sqrt{5}}{2} \approx 0.382
\end{equation}
By always choosing the next point exactly $38.2\%$ into the larger segment, the algorithm guarantees optimal, predictable shrinking of the search space regardless of how malicious the function's shape is.

\subsection{Minimization in a Multidimensional Space}
\label{subsec:multimin}
When finding a multidimensional minimum, gradient descent leads us to a flat spot ($\nabla f = 0$). But is it a minimum, a maximum, or a saddle point (like a mountain pass)?
To confirm a true minimum, the landscape must curve upwards in \textit{every} direction. Mathematically, the \textbf{Hessian Matrix} of second derivatives ($A_{ij} = \frac{\partial^2 f}{\partial x_i \partial x_j}$) must be \textbf{Positive Definite}.

Calculating eigenvalues to prove this is computationally expensive. A beautiful trick is to apply Gaussian elimination (see Section~\ref{subsec:lineareq}) to the Hessian. If and only if every single pivot element on the diagonal ($B_{ii}$) remains strictly positive during elimination, the matrix is guaranteed to be positive definite, confirming we have found the true bottom of the valley.

\vspace{0.5cm}
\noindent \textbf{Further Reading for Lecture \ref{sec:lec4}:}
\begin{itemize}
    \item \textit{Numerical Recipes}, W.H. Press et al. - Chapter 11.1 (Jacobi Transformations), Chapter 9.7 (Globally Convergent Methods for Nonlinear Systems), Chapter 10.1 (Golden Section Search).
    \item \textit{Computational Physics}, Mark Newman - Chapter 6 (Solution of Linear and Nonlinear Equations).
\end{itemize}

\end{document}
